\documentclass[11pt, reqno]{amsart}

\usepackage[spanish]{babel}
\usepackage[LGR, T1]{fontenc}
\usepackage[utf8]{inputenc}

../LaTeX/general.tex
% \input{../graphics.tex}

\makeatletter
\def\emailaddrname{\textit{Correo electrónico}}
\def\subtitle#1{\gdef\@subtitle{#1}}
\def\@subtitle{}

% Metadata
\def\logo#1{\gdef\@logo{#1}}
\def\@logo{}
\def\institution#1{\gdef\@institution{#1}}
\def\@institution{}
\def\department#1{\gdef\@department{#1}}
\def\@department{}
\def\professor#1{\gdef\@professor{#1}}
\def\@professor{}
\def\course#1{\gdef\@course{#1}}
\def\@course{}
\def\coursecode#1{\gdef\@coursecode{#1}}
\def\@coursecode{}

\renewcommand{\maketitle}{
\begin{center}
	\small
	\renewcommand{\arraystretch}{1.2}
	\begin{tabular}{cp{.37\textwidth}p{0.44\textwidth}}
		% \hline
		\multirow{5}{*}{\includegraphics[height=2.0cm]{\@logo}}
	  & \multicolumn{2}{c}{ \makecell{{\bfseries \@institution} \\ \@department} } \\
	  % & \multicolumn{2}{c|}{{\bfseries\@institution} \\ \@department} \\
	  \cline{2-3}
	  & \textbf{Profesor:} \@professor & \textbf{Ayudante:} \authors \\
	  % \cline{2-3}
	  & \textbf{Curso:} \@course & \textbf{Sigla:} \@coursecode \\
	  % \cline{2-3}
	  & \multicolumn{2}{l}{ \textbf{Fecha:} \@date } \\
	  % \hline
	\end{tabular}
	\\[\baselineskip]
	% {}
	% \vspace{2\baselineskip}
	{\bfseries\Large\@title}
	\ifx\@subtitle\@empty\else
		\\[1ex]
		\large\mdseries\@subtitle
	\fi
\end{center}
}
\makeatother

\usepackage{multirow, makecell}

\usepackage[
	reversemp,
	letterpaper,
	% marginpar=2cm,
	% marginsep=1pt,
	margin=2.3cm
]{geometry}
\usepackage{fontawesome}
% \makeatletter
% \@reversemargintrue
% \makeatother

% Símbolos al margen, necesitan doble compilación
\newcommand{\hard}{\marginnote{\faFire}}
\newcommand{\hhard}{\marginnote{\faFire\faFire}}

% Dependencias para los teoremas
\usepackage{xifthen}
\def\@thmdep{}
\newcommand{\thmdep}[1]{
	\ifthenelse{\isempty{#1}}
	{\def\@thmdep{}}
	{\def\@thmdep{ (#1)}}
}
\newcommand{\thmstyle}{\color{thm}\sffamily\bfseries}

% ===== Estilos de Teoremas ==========
\newtheoremstyle{axiomstyle}
	{0.3cm}
	{0.3cm}
	{\normalfont}
	{0.5cm}
	{\bfseries\scshape}
	{:}
	{4pt}
	{\thmname{#1}\thmnote{ #3}\thmnumber{ (#2)}}
\newtheoremstyle{styleC}
	{0.5cm}
	{0.5cm}
	{\normalfont}
	{0.5cm}
	{\bfseries}
	{:}
	{4pt}
	{\thmname{#1\textrm{\@thmdep}}\thmnumber{ #2}\thmnote{ (#3)}}

% ====== Teoremas (sin borde) ===========
\theoremstyle{axiomstyle}
\newtheorem*{axiom}{Axioma}

% ====== Teoremas (sin borde) ==================
\theoremstyle{styleC}
\newtheorem{thm}{Teorema}[section]
\newtheorem{mydef}[thm]{Definición}
\newtheorem{prop}[thm]{Proposición}
\newtheorem{cor}[thm]{Corolario}
\newtheorem{lem}[thm]{Lema}
\newtheorem{con}[thm]{Conjetura}
\newtheorem*{sol}{Solución}
\newtheorem*{obs}{Observación}
\newtheorem*{ex}{Ejemplo}

\newtcbox{bluebox}[1][]{enhanced jigsaw, 
  sharp corners,
  frame hidden,
  nobeforeafter,
  listing only,
  #1} % comando para crear cajas de colores

\expandafter\let\expandafter\oldproof\csname\string\proof\endcsname
\let\oldendproof\endproof
\renewenvironment{proof}[1][\proofname]{%
  \oldproof[\scshape Demostración:]%
}{\oldendproof} % comando para redefinir la caja de la demostración
\newenvironment{hint}[1][\proofname]{%
  \oldproof[\scshape Pista:]%
}{\oldendproof} % comando para redefinir la caja de la demostración

% colores utilizados
\definecolor{numchap}{RGB}{249,133,29}
\definecolor{chap}{RGB}{6,129,204}
\definecolor{sec}{RGB}{204,0,0}
\definecolor{thm}{RGB}{106,176,240}
\definecolor{thmB}{RGB}{32,31,31}
\definecolor{part}{RGB}{212,66,66}

% ====== Diseño de los titulares ===============
\usepackage[explicit]{titlesec} % para personalizar el documento, la opción <<explicit>> hace que el texto de los titulares sea un objeto interactuable

\titleformat{\subsection}[runin]
	{\bfseries}
	{\textrm{\S}\thesubsection}
	{1ex}
	{#1.}

\setlist[enumerate,1]{label=\arabic*., ref=\arabic*} % Enumerate standards

\usepackage{tikz}
\usetikzlibrary{babel,cd}
% \DeclareMathOperator\Gr{Gr}

\DeclareMathOperator\pr{pr}

\title{El grupo fundamental de Poincaré}
\date{\DTMdate{2026-10-08}}

\author{José Cuevas Barrientos}
\email{josecuevasbtos@uc.cl}
\urladdr{https://josecuevas.xyz/teach/2026-2-ayud/}

\logo{../puc_negro.png}
\institution{Pontificia Universidad Católica de Chile}
\department{Facultad de Matemáticas}
\course{Topología Algebraica}
\coursecode{MAT2850}
\professor{Mauricio Bustamante}

\begin{document}

\maketitle

\section{Recubrimientos}
\begin{enumerate}
	\item\lookright
		Sea $a \colon G\times X \to X$ la acción de un grupo $G$ sobre un espacio topológico $X$ por homeomorfismos
		(i.e., para cada $g \in G$, la aplicación $a(g, -) \colon X \to X$ es un homeomorfismo).
		Se dice que $a$ es \strong{libre y propiamente discontinua} si cada punto $x \in X$ posee un entorno $U$ tal que
		\[
			gU \cap U \ne \emptyset \implies g = 1.
		\]
		% En particular, para nosotros toda acción propiamente discontinua es fiel.
		% \begin{enumerate}
		% 	\item Dada cualquier acción $G \acts X$, 
		% \end{enumerate}
		\par
		Pruebe que si $G \acts X$ es una acción libre y propiamente discontinua, entonces la aplicación cociente $X \epicto G
		\coquot X$ es un recubrimiento.

	\item Pruebe que el recubrimiento de una variedad topológica (con o sin borde) es también una variedad topológica (con o sin borde)
		de la misma dimensión.

	\item Sea $X$ un espacio topológico compacto y sea $f\colon Y \to X$ un recubrimiento.
		Pruebe que $Y$ es compacto syss $f$ tiene fibras finitas (i.e., cada punto posee finitas preimágenes).

		(Para la implicancia \textquote{$\implies$} puede suponer que $X$ es metrizable; no obstante,
		recomendamos al lector más experimentado intentar el problema sin hipótesis adicionales.)
\end{enumerate}

\section{$\pi_1$ y $H_n$}
\begin{enumerate}[resume]
	\item Sea $G$ un grupo.
		Decimos que el espacio topológico $(G, \mathcal{T})$ es un \strong{grupo topológico} si las funciones de inversa y de
		multiplicación son continuas respecto a la topología $\mathcal{T}$.
		\begin{enumerate}
			\item Pruebe que $\pi_1(G, 1)$ es un grupo abeliano.
			\item Concluya que $\pi_1(\Sn[1], *) \cong \Z$.
		\end{enumerate}

	\item\lookst
		Sea $X$ un espacio topológico arbitrario y sea $\overline{\gamma} \in H_2^{\rm sing}(X)$ una clase de homología singular.
		Demuestre que existen funciones continuas $f_j\colon \Sigma_{g_j} \to X$ (donde $\Sigma_g$ es la \textquote{esfera con $g$
		agarraderas} con $g \ge 0$) tal que la suma de las imágenes de los generadores de $H_2(\Sigma_g) \cong \Z$ es
		\[
			\overline{\gamma} = \sum_j H_2(f_j)(1).
		\]
\end{enumerate}

\appendix
\section{Problemas adicionales}
\begin{enumerate}[resume]
	\item Ocasionalmente se dice que una acción $a \colon G \acts X$ por homeomorfismos es \emph{propiamente discontinua} cuando todo
		punto $x$ posee un abierto $U$ para el cuál el conjunto
		\[
			\{ g \in G : gU \cap U \ne \emptyset \}
		\]
		es finito.
		Sea $a \colon G \acts X$ una acción fiel (i.e., todos los estabilizadores son triviales) y propiamente discontinua en este
		sentido.
		\begin{enumerate}
			\item Pruebe que la aplicación cociente $\pi \colon X \to G\coquot X$ siempre es un recubrimiento cuando
				$X$ es un \underline{espacio métrico}.
			\item Dé un contraejemplo de una acción $a$ como antes para la cual la aplicación cociente $\pi \colon X \to
				G\coquot X$ no es un recubrimiento.
			\item ¿Qué condición ha de imponerle a $X$ para que $\pi$ siempre sea un recubrimiento?
				\textcolor{white}{Resp.: $T_3$.}
		\end{enumerate}

	\item Sea $G$ un grupo cualquiera.
		Una \strong{presentación} de $G$, denotada $G = \langle S : R \rangle$ consiste en un complejo de homomorfismos
		\[\begin{tikzcd}
			F(R) \rar["\varphi"] & F(S) \rar["\pi"] & G
			% \rar & 1
		\end{tikzcd}\]
		de modo que $G \cong F(S)/\langle\!\langle \Img \varphi \rangle\!\rangle$, donde $\langle\!\langle T \rangle\!\rangle$
		denota el mínimo subgrupo \emph{normal} que contiene al conjunto $T$.
		% Cuando $R = \langle r_1, \dots, r_n \rangle $, típicamente denotamos a la presentación por $\langle S : r_1 = \cdots = r_n =
		% 1 \rangle$.
		\begin{enumerate}
			\item Muestre que si $G = \langle S : R \rangle$ es una presentación, entonces $\abcl G = \langle S : R \cup \{ [s,
				t] : s, t \in S \} \rangle$, donde $[s, t] = sts^{-1}t^{-1}$ es el conmutador.
			\item Considere\footnotemark{} la presentación $A_5 = \langle x, y : x^5 = y^2 = (xy)^3 = 1 \rangle$.
				\footnotetext{Entiéndase, no demuestre que esto es presentación. Si le interesa, vid.\
				\url{https://math.stackexchange.com/a/1999793}.}
				Calcule $\abcl A_5$.
		\end{enumerate}
\end{enumerate}

\printbibliography

\end{document}
